% ==============================================================================
% 1. 基础设置与字体
% ==============================================================================
\usepackage{xeCJK}
\xeCJKsetup{
    PunctStyle=kaiming,                  % 开明式：句末标点始终压缩为半角，腾出更多空间
    CheckSingle=true                     % 避免段末单字独占一行
}
\tolerance=1000                          % 提高断行容忍度，允许稍松的行以避免标点溢出到下一行
\emergencystretch=3em                    % 加大紧急伸缩量
\setCJKmainfont{SimSun}[
    BoldFont = SimHei,
    ItalicFont = KaiTi
]
\setmainfont{Times New Roman}

% 全局段落设置：无缩进，无段间距，正文行距
\setlength{\parindent}{0pt}
\setlength{\parskip}{0pt}
\renewcommand{\baselinestretch}{1.5}

% 单字不成行、单行不成页
\makeatletter
\newcommand{\examwidoworphanrules}{%
    \clubpenalty=10000%
    \widowpenalty=10000%
    \displaywidowpenalty=10000%
    \brokenpenalty=10000%
    \@clubpenalty=10000%
}
\makeatother
\examwidoworphanrules
\AtBeginDocument{\examwidoworphanrules}
\raggedbottom

% ==============================================================================
% 2. 章节标题样式 (试卷结构)
% ==============================================================================
\ctexset {
    chapter = {
        pagestyle = headings,
        break = \clearpage,      % 关键：每一章(每一份试卷)强制换页
        numbering = false,       % 关键：不显示“第一章”字样，仅显示标题
        format = \bfseries\zihao{3}\centering,  % 三号字
        beforeskip = 0em,
        afterskip = 0em,
    },
    section = {
        format = \bfseries\zihao{-4}, % 小四号字并加粗
        number = \chinese{section}、,  % 格式：一、
        indent = 0pt,
        aftername = \hspace{0pt},
        beforeskip = 0.5em,
        afterskip = 0.5em,
    },
}

\usepackage{etoolbox}
\pretocmd{\chapter}{\problembodyend\setcounter{probnum}{0}\setcounter{section}{0}}{}{}
\pretocmd{\section}{\problembodyend}{}{}
\makeatletter
\patchcmd{\tableofcontents}
  {\@starttoc{toc}}
  {%
    \begingroup
      \normalfont
      \zihao{-4}\selectfont
      \setlength{\columnsep}{1.8em}%
      \setlength{\columnseprule}{0.8pt}%
      \begin{multicols*}{2}\@starttoc{toc}\end{multicols*}%
    \endgroup
  }
  {}{}
\renewcommand*\l@part[2]{%
  \addpenalty{-\@highpenalty}%
  \vskip 0.8em \@plus\p@
  \begingroup
    \parindent \z@
    \rightskip \z@
    \parfillskip \z@
    \centering\normalfont #1\par
  \endgroup
}
\renewcommand*\l@chapter[2]{%
  \ifnum \c@tocdepth >\m@ne
    \addpenalty{\@secpenalty}%
    \addvspace{0pt}%
    \begingroup
      \parindent \z@
      \rightskip \@pnumwidth
      \parfillskip -\@pnumwidth
      \leavevmode
      #1\nobreak\leaders\hbox{$\m@th \mkern 2mu.\mkern 2mu$}\hfill\nobreak
      \hb@xt@\@pnumwidth{\hss #2}\par
    \endgroup
  \fi
}
\makeatother

% 正文小四号字
\AtBeginDocument{\zihao{-4}}

\newcommand{\examspecialsection}[1]{%
    \problembodyend
    \par
    \medskip
    \begin{center}
    \bfseries\zihao{-4}#1
    \end{center}
    \setcounter{section}{0}%
    \medskip
}
\newcommand{\additionalproblemsection}{\examspecialsection{附加题}}
\newcommand{\referenceproblemsection}{\examspecialsection{参考题}}
\newcommand{\optionalproblemsection}{\examspecialsection{选作题}}

% ==============================================================================
% 3. 宏包加载 (精简版)
% ==============================================================================
\usepackage{geometry}
\usepackage{needspace}
\geometry{a4paper,left=2cm,right=2cm,top=2cm,bottom=2cm}
\newcommand{\examneedlines}[1][2]{\Needspace{#1\baselineskip}}
\newif\ifexampageflexused
\AddToHook{shipout/after}{\global\exampageflexusedfalse}
\newcommand{\examavoidshortpage}{%
    \ifexampageflexused\else
        \global\exampageflexusedtrue
        \relax
    \fi
}

% 单栏排版；目录页局部使用 multicol
\usepackage{multicol}

% 数学相关
\usepackage{amsmath, amssymb, amsthm, bm} 
\usepackage{unicode-math}                 
\setmathfont{TeX Gyre Termes Math}
% 希腊字母使用 STIX 字体 (MathType 风格，正体)
\setmathfont{STIX}[
    range = {up/{Greek,greek}, it/{Greek,greek}, bfup/{Greek,greek}, bfit/{Greek,greek}}
]

% 大运算符和空集使用 Latin Modern Math (LaTeX 默认风格)
\setmathfont{Latin Modern Math}[
    range = {\sum, \prod, \coprod, \int, \iint, \iiint, \iiiint, \oint, \bigcup, \bigcap, \bigsqcup, \bigvee, \bigwedge, \bigoplus, \bigotimes, \bigodot}
]
\let\emptyset\varnothing

\setmathfont{Asana Math}[
    range = {\varnothing}
]

\allowdisplaybreaks                       

% cases 行距与正文节奏保持一致；amsmath 的 cases 内部会重设 \arraystretch，
% 所以直接重定义内部入口，避免外层 hook 被覆盖。
\makeatletter
\def\env@cases{%
  \let\@ifnextchar\new@ifnextchar
  \left\lbrace
  \def\arraystretch{0.9}%
  \array{@{}l@{\quad}l@{}}%
}
\makeatother

% 图形与表格
\usepackage{graphicx}
\usepackage{adjustbox}
\usepackage{xparse}
\let\examoriginalsection\section
\RenewDocumentCommand{\section}{s o m}{%
    \IfBooleanTF{#1}{%
        \IfValueTF{#2}{\examoriginalsection*[#2]{#3}}{\examoriginalsection*{#3}}%
    }{%
        \ifstrequal{#3}{附加题}{%
            \additionalproblemsection
        }{%
            \ifstrequal{#3}{参考题}{%
                \referenceproblemsection
            }{%
                \ifstrequal{#3}{选作题}{%
                    \optionalproblemsection
                }{%
                    \IfValueTF{#2}{\examoriginalsection[#2]{#3}}{\examoriginalsection{#3}}%
                }%
            }%
        }%
    }%
}
\usepackage{tikz}                         
\usetikzlibrary{arrows.meta, calc, intersections, positioning, shapes.geometric}
\usepackage{wrapfig}
% 调整 wrapfigure 与正文的间距
\setlength{\intextsep}{0pt}  % 纵向：图片上下的间距
% \columnsep（横向：图片与侧边文字的间距）已在上方 multicol 区域统一设为 1.5em
\usepackage{array, tabularray, multirow}  
\usepackage{caption}

% --- 接管 \includegraphics：找不到文件时显示占位框而非报错 ---
\let\originalincludegraphics\includegraphics
\renewcommand{\includegraphics}[2][]{%
    \problemconsumeinlineblockprefix
    \IfFileExists{#2}{%
        \originalincludegraphics[#1]{#2}%
    }{%
        \setlength{\fboxsep}{4pt}%
        \fbox{%
            \begin{minipage}[c][4em][c]{0.8\linewidth}%
                \centering\color{gray}\small 图片未找到：{\ttfamily\detokenize{#2}}%
            \end{minipage}%
        }%
    }%
}

\ExplSyntaxOn
\bool_new:N \l__bitmap_solution_caption_bool
\tl_new:N \l__bitmap_solution_caption_tl

\cs_new_protected:Npn \__bitmap_prepare_caption:n #1
{
    \bool_set_false:N \l__bitmap_solution_caption_bool
    \tl_clear:N \l__bitmap_solution_caption_tl
    \regex_extract_once:nnNT { q([0-9]+)_solution_fig([0-9]+)\.(?:tex|png|jpg|jpeg)$ } {#1} \l_tmpa_seq
    {
        \bool_set_true:N \l__bitmap_solution_caption_bool
        \tl_set:Nx \l__bitmap_solution_caption_tl
        {第\seq_item:Nn \l_tmpa_seq {2}题解析图\seq_item:Nn \l_tmpa_seq {3}}
    }
}
\NewDocumentCommand{\bitmappreparecaption}{m}
{
    \__bitmap_prepare_caption:n {#1}
}
\NewDocumentCommand{\bitmapinsertcaption}{}
{
    \bool_if:NTF \l__bitmap_solution_caption_bool
        {\captionof*{figure}{\tl_use:N \l__bitmap_solution_caption_tl}}
        {\captionof{figure}{}}
}
\NewDocumentCommand{\bitmapcaptiontext}{m}
{
    \bitmappreparecaption{#1}
    \bool_if:NTF \l__bitmap_solution_caption_bool
        {\tl_use:N \l__bitmap_solution_caption_tl}
        {第\theprobnum 题图}
}
\ExplSyntaxOff

\newcommand{\bitmapinclude}[2][]{%
    \begingroup
    \adjustbox{max width=\linewidth,max totalheight=4.8cm,#1}{\includegraphics{#2}}%
    \endgroup
}

\newcommand{\bitmapfigure}[2][]{%
    \begingroup
    \bitmappreparecaption{#2}%
    \begin{center}
    \bitmapinclude[#1]{#2}
    \bitmapinsertcaption
    \end{center}
    \endgroup
}

\NewDocumentCommand{\bitmapinlinefigure}{O{0.42\linewidth} O{} m}
{%
    \begingroup
    \bitmappreparecaption{#3}%
    \makebox[#1][c]{%
        \vtop{%
            \hsize=#1\relax
            \centering
            \bitmapinclude[#2]{#3}\par
            \small \bitmapcaptiontext{#3}\par
        }%
    }%
    \endgroup
}

% 列表与枚举
\usepackage{enumitem}
\AtBeginEnvironment{center}{\problemcenterstart}
\AfterEndEnvironment{center}{\problemcenterend}
\AtBeginEnvironment{tabular}{\problemconsumeinlineblockprefix}

% 盒子宏包 (用于答案)
\usepackage[most]{tcolorbox}

\usepackage{ulem}
\newcommand{\fillinblank}{\mbox{\ttfamily\_\_\_\_\_\_\_\_}}

% ==============================================================================
% 4. 核心命令：题目、选项
% ==============================================================================

% --- A. 题目计数器 ---
% 绑定到 chapter，每次新试卷(chapter)自动重置为 0
% 注意：虽然 chapter 设置了 numbering=false，但计数器依然在后台运作
\newcounter{probnum}[chapter]
\newlength{\problembodyindent}
\setlength{\problembodyindent}{1.5em}
\newif\ifproblembodyactive
\newif\ifproblemfirstparagraph
\newif\ifproblemheadpending

\newcommand{\problemresetlayout}{%
    \leftskip=0pt\relax
    \hangindent=0pt\relax
    \hangafter=1\relax
    \parshape=0\relax
}

\newcommand{\problemheadbox}{%
    \makebox[\problembodyindent][l]{\textbf{\theprobnum.}}%
}

\newcommand{\problemheadprefix}{%
    \textbf{\theprobnum.}\hspace{0.5em}%
}
\newcommand{\problemoriginallabel}{}%

\newcommand{\problembodyfinish}{%
    \ifhmode\par\fi
    \global\problembodyactivefalse
    \global\problemfirstparagraphfalse
    \global\problemheadpendingfalse
    \everypar{}%
    \problemresetlayout
}
\let\problembodyend\problembodyfinish

\newcommand{\problemliststart}{%
    \ifproblembodyactive
        \ifhmode\unskip\par\fi
        \ifproblemheadpending
            \global\problemfirstparagraphfalse
            \let\problemoriginallabel\makelabel
            \def\makelabel##1{%
                \ifproblemheadpending
                    \global\problemheadpendingfalse
                    \llap{\problemheadprefix}%
                \fi
                \problemoriginallabel{##1}%
            }%
        \fi
    \fi
}

\newcommand{\problemblockstart}{%
    \ifproblembodyactive
        \ifhmode\unskip\par\fi
        \ifproblemheadpending
            \global\problemheadpendingfalse
            \global\problemfirstparagraphfalse
            \llap{\problemheadprefix}%
        \fi
    \fi
}

\newcommand{\problemconsumeinlineblockprefix}{%
    \ifproblemheadpending
        \global\problemheadpendingfalse
        \leavevmode\llap{\problemheadprefix}%
    \fi
}

\newcommand{\problemcenterstart}{%
    \ifproblembodyactive
        \ifhmode\unskip\par\fi
        \ifproblemheadpending
            \global\problemfirstparagraphfalse
            \linewidth=\dimexpr\linewidth-\problembodyindent\relax
            \hsize=\linewidth
        \fi
    \fi
}

\newcommand{\problemcenterend}{%
    \ifproblembodyactive
        \global\everypar{\problemcontinuationpar}%
    \fi
}

\makeatletter
\patchcmd{\@endparenv}
  {\addpenalty\@endparpenalty\addvspace\@topsepadd\@endpetrue}
  {\addpenalty\@endparpenalty\addvspace\@topsepadd%
   \ifproblembodyactive
     \global\@endpefalse
     \global\everypar{\problemcontinuationpar}%
   \else
     \@endpetrue
   \fi}
  {}{}
\makeatother

\newcommand{\problemcontinuationpar}{%
    \ifproblembodyactive
        \ifproblemfirstparagraph
            \leftskip=0pt\relax
            \hangindent=\problembodyindent
            \hangafter=1
            \parshape=0\relax
            \ifproblemheadpending
                \problemheadbox
                \global\problemheadpendingfalse
            \fi
            \global\problemfirstparagraphfalse
        \else
            \leftskip=\problembodyindent\relax
            \hangindent=0pt\relax
            \hangafter=1\relax
            \parshape=0\relax
        \fi
    \fi
}

% --- B. 题目命令 \problem ---
\newcommand{\problem}{%
    \problembodyfinish
    \examavoidshortpage
    \examneedlines[1]%
    \stepcounter{probnum}%
    \global\problembodyactivetrue
    \global\problemfirstparagraphtrue
    \global\problemheadpendingtrue
    \everypar{\problemcontinuationpar}%
    \ignorespaces%
}
\let\endproblem\problembodyfinish

% --- C. 智能选项 \choices ---
\newcommand{\choiceentry}[2]{%
    \hangindent=1.8em\hangafter=1\noindent\textbf{#1.}~#2%
}

\newcommand{\choicecell}[3]{%
    \parbox[t]{#1}{\raggedright\textbf{#2.}~#3}%
}

\newdimen\examchoicewidth
\newdimen\examchoicewidthA
\newdimen\examchoicewidthB
\newdimen\examchoicewidthC
\newdimen\examchoicewidthD
\newdimen\examchoicemaxwidth

\newcommand{\choices}[4]{%
    \problembodyend
    \par\nopagebreak
    \examneedlines[2]%
    \vspace{4pt}
    \begingroup
    \setlength{\leftskip}{\problembodyindent}
    \examchoicewidth=\dimexpr\linewidth-\problembodyindent\relax
    \settowidth{\examchoicewidthA}{\textbf{A.}~#1}
    \settowidth{\examchoicewidthB}{\textbf{B.}~#2}
    \settowidth{\examchoicewidthC}{\textbf{C.}~#3}
    \settowidth{\examchoicewidthD}{\textbf{D.}~#4}
    \examchoicemaxwidth=\examchoicewidthA
    \ifdim\examchoicewidthB>\examchoicemaxwidth \examchoicemaxwidth=\examchoicewidthB \fi
    \ifdim\examchoicewidthC>\examchoicemaxwidth \examchoicemaxwidth=\examchoicewidthC \fi
    \ifdim\examchoicewidthD>\examchoicemaxwidth \examchoicemaxwidth=\examchoicewidthD \fi
    \addtolength{\examchoicemaxwidth}{2em}

    \vspace{-0.2em}\noindent
    \ifdim\examchoicemaxwidth<0.22\examchoicewidth
        \choicecell{\dimexpr0.22\examchoicewidth\relax}{A}{#1}\hfill
        \choicecell{\dimexpr0.22\examchoicewidth\relax}{B}{#2}\hfill
        \choicecell{\dimexpr0.22\examchoicewidth\relax}{C}{#3}\hfill
        \choicecell{\dimexpr0.22\examchoicewidth\relax}{D}{#4}%
    \else
        \ifdim\examchoicemaxwidth<0.47\examchoicewidth
            \choicecell{\dimexpr0.47\examchoicewidth\relax}{A}{#1}\hfill
            \choicecell{\dimexpr0.47\examchoicewidth\relax}{B}{#2}\par
            \vspace{0.2em}
            \choicecell{\dimexpr0.47\examchoicewidth\relax}{C}{#3}\hfill
            \choicecell{\dimexpr0.47\examchoicewidth\relax}{D}{#4}%
        \else
            \choiceentry{A}{#1}\par
            \choiceentry{B}{#2}\par
            \choiceentry{C}{#3}\par
            \choiceentry{D}{#4}%
        \fi
    \fi
    \par
    \endgroup
}

% ==============================================================================
% 5. 答案控制系统 (框式设计)
% ==============================================================================
\newif\ifshowanswer 
\showanswertrue     % 【默认开启】

% 定义答案盒子的样式
\tcbset{
    examanswer/.style={
        enhanced jigsaw,
        breakable,             % 允许跨栏
        colback=gray!5,        % 极淡的灰色背景
        colframe=gray!40,      % 边框颜色
        boxrule=0.5pt,         % 细边框
        arc=1mm,               % 轻微圆角
        left=3pt,right=2pt, top=1pt, bottom=1pt,
        before skip=0.3em,
        after skip=0.3em,
        before upper={\examwidoworphanrules\setlength{\parindent}{2em}\tolerance=3200\emergencystretch=8em\setlength{\jot}{4pt}\lineskiplimit=3pt\lineskip=4pt\selectfont},
        fontupper=\zihao{-4},
    }
}

% 答案/解析环境；旧命令写法 \answer{...}、\solution{...} 继续兼容。
\newbox\examhiddenanswerbox
\newcommand{\answerbegin}{%
    \problembodyend
    \ifshowanswer
        \examneedlines[3]%
        \begin{tcolorbox}[examanswer]
            \textbf{【答案】}\ignorespaces
    \else
        \setbox\examhiddenanswerbox=\vbox\bgroup\hsize=\linewidth
    \fi
}
\newcommand{\answerend}{%
    \ifshowanswer
        \end{tcolorbox}
    \else
        \egroup
    \fi
}
\NewDocumentCommand{\answer}{+g}{%
    \IfNoValueTF{#1}{\answerbegin}{\answerbegin #1\answerend}%
}
\let\endanswer\answerend

\newcommand{\solutionbegin}{%
    \problembodyend
    \ifshowanswer
        \examneedlines[3]%
        \begin{tcolorbox}[examanswer]
            \textbf{【解析】}\ignorespaces
    \else
        \setbox\examhiddenanswerbox=\vbox\bgroup\hsize=\linewidth
    \fi
}
\newcommand{\solutionend}{%
    \ifshowanswer
        \end{tcolorbox}
    \else
        \egroup
    \fi
}
\NewDocumentCommand{\solution}{+g}{%
    \IfNoValueTF{#1}{\solutionbegin}{\solutionbegin #1\solutionend}%
}
\let\endsolution\solutionend

% ==============================================================================
% 5.1 解析区行间公式归一化宏
% ==============================================================================
\NewDocumentCommand{\examdisplaymath}{+m}{%
    \[
    #1
    \]
}
\NewDocumentCommand{\examdisplayaligned}{o +m}{%
    \[
    \IfValueTF{#1}{\begin{aligned}[#1]}{\begin{aligned}}
    #2
    \end{aligned}
    \]
}
\NewDocumentCommand{\examdisplaycases}{+m +m}{%
    \[
    #1\begin{cases}
    #2
    \end{cases}
    \]
}
\NewDocumentCommand{\examdisplayarray}{+m m +m +m}{%
    \[
    #1\begin{array}{#2}
    #3
    \end{array}#4
    \]
}
\NewDocumentCommand{\examdisplaychain}{+m}{%
    \begin{align*}
    #1
    \end{align*}
}

% ==============================================================================
% 6. 图表标题定制 (第X题图)
% ==============================================================================
\DeclareCaptionLabelFormat{examfig}{第\theprobnum 题图}
\DeclareCaptionLabelFormat{examtab}{第\theprobnum 题表}

\captionsetup[figure]{
    labelformat=examfig, 
    labelsep=none,       
    font=small
}
\captionsetup[table]{
    labelformat=examtab, 
    labelsep=none,
    font=small
}

% ==============================================================================
% 7. 辅助符号与工具
% ==============================================================================
\renewcommand{\thefootnote}{\circled{\arabic{footnote}}}
\newcommand{\circled}[2][]{\tikz[baseline=(char.base)]
    {\node[shape=circle,draw,inner sep=0.8pt](char){\small #2};}}

\newcommand{\R}{\symbfup{R}}
\newcommand{\Z}{\symbfup{Z}}
\newcommand{\N}{\mathbb{N}}
\renewcommand{\Pr}{P}
\renewcommand{\d}{\text{d}}
\renewcommand{\i}{\text{i}}
\newcommand{\e}{\text{e}}
\newcommand{\degree}{^\circ}
\newcommand{\bs}[1]{\symbfit{#1}} 
\let\oldfrac\frac
\renewcommand{\frac}[2]{%
    \mathchoice
        {\dfrac{#1}{#2}}                             % 1. Display style (行间公式时) -> 使用大分式
        {\dfrac{#1}{#2}}                             % 2. Text style (行内公式时) -> 强制使用大分式
        {\oldfrac{\scriptstyle #1}{\scriptstyle #2}} % 3. Script style (指数、上下标) -> 强制分子分母维持 scriptstyle
        {\oldfrac{\scriptscriptstyle #1}{\scriptscriptstyle #2}} % 4. Scriptscript style (更深层的下标) -> 恢复默认小分式
}

% 大运算符在行内公式中默认使用显示模式的格式（类似 \frac 的处理）
\makeatletter
% 辅助宏：sum 类大运算符（\mathchoice 控制大小，\limits 控制极限位置）
\newcommand{\make@display@bigop}[2]{% #1=备份命令, #2=原命令
  \renewcommand{#2}{\mathop{%
    \mathchoice
        {#1}%                              % 1. Display style -> 保持原样（已经是大符号）
        {\displaystyle #1}%                % 2. Text style (行内公式) -> 强制大符号
        {#1}%                              % 3. Script style (指数、上下标) -> 保持原样
        {#1}%                              % 4. Scriptscript style -> 保持原样
  }\limits}}
% 辅助宏：积分类大运算符（同理，但用 \nolimits 保持极限在侧边）
\newcommand{\make@display@bigint}[2]{% #1=备份命令, #2=原命令
  \renewcommand{#2}{\mathop{%
    \mathchoice
        {#1}%                              % 1. Display style -> 保持原样
        {\displaystyle #1}%                % 2. Text style (行内公式) -> 强制大符号
        {#1}%                              % 3. Script style -> 保持原样
        {#1}%                              % 4. Scriptscript style -> 保持原样
  }\nolimits}}
\makeatother

% \let\oldlim\lim
% \renewcommand{\lim}{\oldlim\limits}

\newcommand{\myarc}[1]{\overset{\frown}{#1}}

\SetEnumitemKey{examhang}{%
    nosep,
    align=left,
    before=\problemliststart,
    itemindent=0pt,
    listparindent=0pt,
    parsep=0pt,
    topsep=0pt,
    partopsep=0pt
}

\setlist[enumerate,1]{ % 仅针对第一层级
    examhang,
    label=（\arabic*）,
    left=\problembodyindent,
    labelsep=0.5em,
    widest*=9
}
\setlist[enumerate,2]{
    examhang,
    label=（\roman*）,
    % left=2.5em\relax,
    labelsep=0.5em,
    widest*=8
}

\newlist{circlelist}{enumerate}{1}
\setlist[circlelist,1]{
    examhang,
    label=\circled{\arabic*},
    left=\problembodyindent,
    labelsep=0.5em,
    widest*=9
}

% 修改公式与前后文的紧凑度
\makeatletter
\g@addto@macro\normalsize{
    % --- 控制行间公式 (Display Math) ---
	\setlength\abovedisplayskip{0.25em plus 0.5em minus 0.25em}
	\setlength\belowdisplayskip{0.25em plus 0.5em minus 0.25em}
	\setlength\abovedisplayshortskip{0em plus 2pt minus 0pt}
	\setlength\belowdisplayshortskip{0em plus 2pt minus 0pt}

    % --- 控制行内公式 (Inline Math) 的避让机制 ---
    % 设置阈值：当两行间距小于 1pt 时触发避让
    \setlength\lineskiplimit{3pt} 
    % 设置避让间距：触发后强制留出 3pt 的缝隙，防止重叠
    \setlength\lineskip{3pt}
}
\makeatother

\makeatletter
\AtBeginDocument{
    % 处理 pi 为正体
    \let\oldpi\pi
    \renewcommand{\pi}{\symup\oldpi}

    % 组合数统一改成高中记号 C_n^m；无论源码使用 \binom / \dbinom / \tbinom，
    % 成稿都输出同一套样式。
    \RenewDocumentCommand{\binom}{m m}{\mathrm C_{#1}^{#2}}
    \RenewDocumentCommand{\dbinom}{m m}{\mathrm C_{#1}^{#2}}
    \RenewDocumentCommand{\tbinom}{m m}{\mathrm C_{#1}^{#2}}
    
    % 处理 nabla：不带 ^2 时加粗，带 ^2 时（Laplacian）保持原样
    \let\oldnabla\nabla
    \RenewDocumentCommand{\nabla}{e{^}}{%
        \IfValueTF{#1}{% 有上标
            \ifstrequal{#1}{2}{% 上标是 2
                \oldnabla^{#1}%
            }{% 上标不是 2
                \symbf{\oldnabla}^{#1}%
            }%
        }{% 无上标
            \symbf{\oldnabla}%
        }%
    }

    % --- 大运算符：在 \AtBeginDocument 中重定义，确保在 unicode-math 之后 ---
    % sum 类
    \let\oldsum\sum             \make@display@bigop{\oldsum}{\sum}
    \let\oldprod\prod           \make@display@bigop{\oldprod}{\prod}
    \let\oldcoprod\coprod       \make@display@bigop{\oldcoprod}{\coprod}
    \let\oldbigcup\bigcup       \make@display@bigop{\oldbigcup}{\bigcup}
    \let\oldbigcap\bigcap       \make@display@bigop{\oldbigcap}{\bigcap}
    \let\oldbigsqcup\bigsqcup   \make@display@bigop{\oldbigsqcup}{\bigsqcup}
    \let\oldbigvee\bigvee       \make@display@bigop{\oldbigvee}{\bigvee}
    \let\oldbigwedge\bigwedge   \make@display@bigop{\oldbigwedge}{\bigwedge}
    \let\oldbigoplus\bigoplus   \make@display@bigop{\oldbigoplus}{\bigoplus}
    \let\oldbigotimes\bigotimes \make@display@bigop{\oldbigotimes}{\bigotimes}
    \let\oldbigodot\bigodot     \make@display@bigop{\oldbigodot}{\bigodot}
    % 积分类
    \let\oldint\int             \make@display@bigint{\oldint}{\int}
    \let\oldiint\iint           \make@display@bigint{\oldiint}{\iint}
    \let\oldiiint\iiint         \make@display@bigint{\oldiiint}{\iiint}
    \let\oldiiiint\iiiint       \make@display@bigint{\oldiiiint}{\iiiint}
    \let\oldoint\oint           \make@display@bigint{\oldoint}{\oint}
}
\makeatother
